\documentclass[12pt,a4paper]{report}

% Packages
\usepackage[utf8]{inputenc}
\usepackage[T1]{fontenc}
\usepackage[english]{babel}
\usepackage{geometry}
\geometry{margin=1in}

\usepackage{listings}
\usepackage{hyperref}
\usepackage{graphicx}
\usepackage{float}
\usepackage{booktabs}
\usepackage{xcolor}

% Code formatting
\lstset{
    basicstyle=\ttfamily\footnotesize,
    breaklines=true,
    captionpos=b,
    numbers=left,
    numberstyle=\tiny,
    frame=single,
    backgroundcolor=\color{gray!10},
    keywordstyle=\color{blue},
    commentstyle=\color{green!60!black},
    stringstyle=\color{red},
    showstringspaces=false
}

% Hyperref setup
\hypersetup{
    colorlinks=true,
    linkcolor=blue,
    filecolor=magenta,
    urlcolor=cyan,
    citecolor=red
}

% Title page information
\title{Tunisia National Parks API: \\ A Comprehensive Web Service for Biodiversity and Environmental Data Management}
\author{Mohamed Ali Ben Salah}
\date{January 17, 2026}

\begin{document}

% Title page
\maketitle

% Table of contents
\tableofcontents
\newpage

% Introduction
\chapter{Introduction}

\section{Project Overview}

The Tunisia National Parks API represents a comprehensive web service designed to provide structured access to biodiversity and environmental data for Tunisia's protected natural areas. This project addresses the critical need for centralized, programmatically accessible information about national parks, wildlife species, and environmental monitoring data.

\section{Problem Statement}

Tunisia's rich biodiversity and extensive network of national parks lack a unified, modern data management system. Environmental researchers, conservation organizations, and tourism operators face significant challenges in accessing accurate, up-to-date information about park locations, species distributions, weather conditions, and visitor facilities. The absence of a standardized API hinders data-driven decision making in environmental conservation and sustainable tourism development.

\section{Objectives}

The primary objectives of this project include:

\begin{itemize}
    \item Develop a RESTful API providing comprehensive access to Tunisia's national parks data
    \item Implement robust data models for parks, species, trails, and environmental monitoring
    \item Integrate external services for weather, imagery, and geographical data
    \item Create an intuitive web interface for data visualization and management
    \item Ensure scalability, security, and maintainability of the system
\end{itemize}

\section{Scope and Limitations}

The project encompasses 18 national parks and protected areas, covering approximately 250 species of flora and fauna. The system integrates with external APIs for enhanced functionality while maintaining data sovereignty for core environmental information.

% System Architecture
\chapter{System Architecture}

\section{High-Level Architecture}

The Tunisia National Parks API follows a modern web application architecture comprising three primary layers: presentation, application, and data persistence. The system employs a microservices-inspired approach with clear separation of concerns between frontend and backend components.

\begin{figure}[H]
    \centering
    \includegraphics[width=\textwidth]{system_architecture.png}
    \caption{System Architecture Overview}
    \label{fig:architecture}
\end{figure}

\section{Technology Stack}

\subsection{Backend Technologies}

The backend implementation utilizes Python-based technologies optimized for API development and data processing:

\begin{itemize}
    \item \textbf{FastAPI}: High-performance web framework for building REST APIs
    \item \textbf{Python 3.11+}: Core programming language with advanced type hinting
    \item \textbf{SQLAlchemy/SQLModel}: Object-relational mapping for database operations
    \item \textbf{Pydantic}: Data validation and serialization
    \item \textbf{Uvicorn}: ASGI server for production deployment
\end{itemize}

\subsection{Frontend Technologies}

The presentation layer employs modern web technologies for responsive user interfaces:

\begin{itemize}
    \item \textbf{HTML5/CSS3}: Semantic markup and responsive styling
    \item \textbf{Vanilla JavaScript}: Client-side interactivity without framework dependencies
    \item \textbf{Jinja2}: Server-side templating for dynamic content
    \item \textbf{Leaflet.js}: Interactive mapping capabilities
    \item \textbf{Font Awesome}: Icon library for enhanced user experience
\end{itemize}

\subsection{Database Layer}

The data persistence layer supports multiple database systems for flexibility:

\begin{itemize}
    \item \textbf{SQLite}: Development and lightweight deployments
    \item \textbf{PostgreSQL/PostGIS}: Production environments with spatial capabilities
    \item \textbf{Alembic}: Database migration management
\end{itemize}

\section{External Integrations}

The system integrates with several external APIs to enhance functionality:

\begin{itemize}
    \item \textbf{OpenWeatherMap API}: Real-time weather data and forecasts
    \item \textbf{Unsplash API}: High-quality nature photography
    \item \textbf{NewsAPI}: Environmental news and conservation updates
    \item \textbf{OpenStreetMap}: Geographical mapping data
\end{itemize}

% Backend Implementation
\chapter{Backend Implementation}

\section{Application Framework}

The backend is built using FastAPI, a modern, high-performance web framework for building APIs with Python 3.7+ based on standard Python type hints. FastAPI provides automatic API documentation generation, data validation, and asynchronous request handling capabilities.

\subsection{Core Application Structure}

\begin{lstlisting}[language=Python, caption=Main Application Entry Point]
from fastapi import FastAPI, HTTPException, Request
from fastapi.middleware.cors import CORSMiddleware
from fastapi.staticfiles import StaticFiles
from fastapi.templating import Jinja2Templates
from fastapi.responses import HTMLResponse

from database import init_db
from routers import parks, species, auth
from config import settings

# Initialize FastAPI application
app = FastAPI(
    title="Tunisia National Parks API",
    description="Complete API for Tunisia's national parks",
    version="3.0.0",
    docs_url="/docs"
)

# CORS middleware
app.add_middleware(
    CORSMiddleware,
    allow_origins=settings.get_origins_list(),
    allow_credentials=True,
    allow_methods=["*"],
    allow_headers=["*"],
)

# Static files and templates
app.mount("/static", StaticFiles(directory="static"), name="static")
templates = Jinja2Templates(directory="templates")

# Include routers
app.include_router(parks.router)
app.include_router(species.router)
app.include_router(auth.router)
\end{lstlisting}

\section{Database Schema}

The database schema is designed to support complex relationships between geographical, biological, and user-generated data. The system employs a relational database structure optimized for spatial queries and hierarchical data management.

\subsection{Core Tables}

\begin{itemize}
    \item \textbf{parks}: Geographic and administrative information for national parks
    \item \textbf{species}: Comprehensive species database with conservation status
    \item \textbf{park\_species}: Many-to-many relationship between parks and species
    \item \textbf{trails}: Hiking trail information with difficulty ratings
    \item \textbf{reviews}: User-generated park reviews and ratings
    \item \textbf{users}: User account management and authentication
\end{itemize}

\subsection{Data Models}

\begin{lstlisting}[language=Python, caption=Park Data Model]
from sqlmodel import SQLModel, Field, Relationship
from typing import List, Optional
from datetime import datetime

class ParkDB(SQLModel, table=True):
    __tablename__ = "parks"

    id: int = Field(default=None, primary_key=True)
    name: str = Field(index=True)
    governorate: str = Field(index=True)
    description: str
    latitude: float
    longitude: float
    area_km2: float
    images: Optional[List[str]] = Field(default=None, sa_column=Column(JSON))
    google_maps_url: str
    created_at: datetime = Field(default_factory=datetime.utcnow)
    updated_at: datetime = Field(default_factory=datetime.utcnow)

    # Relationships
    species: List["ParkSpeciesLink"] = Relationship(back_populates="park")
    trails: List["TrailDB"] = Relationship(back_populates="park")
    reviews: List["ReviewDB"] = Relationship(back_populates="park")
\end{lstlisting}

\section{API Routing and Controllers}

The application employs a modular routing system with dedicated routers for different functional domains. Each router handles specific business logic and data operations.

\subsection{Parks Router}

\begin{lstlisting}[language=Python, caption=Parks API Endpoints]
from fastapi import APIRouter, HTTPException, Query
from sqlmodel import Session, select
from typing import List

from database import get_engine
from models import ParkDB, Park

router = APIRouter(prefix="/api/parks", tags=["Parks"])

@router.get("/", response_model=List[Park])
def list_parks(
    skip: int = Query(0, ge=0),
    limit: int = Query(50, ge=1, le=100),
    governorate: str | None = None,
):
    """List all parks with advanced filtering and pagination."""
    with Session(get_engine()) as session:
        statement = select(ParkDB)

        if governorate:
            statement = statement.where(ParkDB.governorate == governorate)

        statement = statement.offset(skip).limit(limit)
        parks_db = session.exec(statement).all()

        return [
            Park(
                id=p.id,
                name=p.name,
                governorate=p.governorate,
                description=p.description,
                latitude=p.latitude,
                longitude=p.longitude,
                area_km2=p.area_km2,
                images=[f"/uploads/parks/{p.id}/{img}" for img in (p.images or [])],
            )
            for p in parks_db
        ]
\end{lstlisting}

\section{Security Implementation}

The system implements multiple layers of security to protect sensitive data and prevent unauthorized access.

\subsection{Authentication System}

\begin{lstlisting}[language=Python, caption=JWT Authentication]
from datetime import datetime, timedelta
from jose import JWTError, jwt
from passlib.context import CryptContext

SECRET_KEY = "your-secret-key"
ALGORITHM = "HS256"
ACCESS_TOKEN_EXPIRE_MINUTES = 30

pwd_context = CryptContext(schemes=["pbkdf2_sha256"], deprecated="auto")

def create_access_token(data: dict):
    to_encode = data.copy()
    expire = datetime.utcnow() + timedelta(minutes=ACCESS_TOKEN_EXPIRE_MINUTES)
    to_encode.update({"exp": expire})
    encoded_jwt = jwt.encode(to_encode, SECRET_KEY, algorithm=ALGORITHM)
    return encoded_jwt

def verify_token(token: str):
    try:
        payload = jwt.decode(token, SECRET_KEY, algorithms=[ALGORITHM])
        return payload
    except JWTError:
        return None
\end{lstlisting}

% Frontend Implementation
\chapter{Frontend Implementation}

\section{User Interface Architecture}

The frontend implementation follows a component-based architecture using vanilla JavaScript and modern HTML5/CSS3 standards. The system employs server-side rendering with Jinja2 templates for optimal performance and SEO.

\subsection{Template Structure}

\begin{lstlisting}[language=HTML, caption=Jinja2 Template Example]
{% extends "base.html" %}

{% block title %}{{ park.name }} - Tunisia National Parks{% endblock %}

{% block content %}
<section class="park-header">
    <div class="park-header-overlay"></div>
    <div class="container">
        <div class="park-header-content">
            <h1>{{ park.name }}</h1>
            <div class="park-meta">
                <span><i class="fas fa-map-marker-alt"></i> {{ park.governorate }}</span>
                <span><i class="fas fa-ruler-combined"></i> {{ park.area_km2 }} km²</span>
            </div>
        </div>
    </div>
</section>

<div class="container">
    <div class="park-details">
        <h2>About {{ park.name }}</h2>
        <p>{{ park.description }}</p>

        <!-- Dynamic content loaded via JavaScript -->
        <div id="weatherInfo"></div>
        <div id="speciesList"></div>
        <div id="reviewsSection"></div>
    </div>
</div>
{% endblock %}
\end{lstlisting}

\section{JavaScript Components}

The frontend employs modular JavaScript for enhanced interactivity and API consumption.

\subsection{Weather Component}

\begin{lstlisting}[language=JavaScript, caption=Weather Data Fetching]
async function loadWeather() {
    const parkId = {{ park.id }};

    try {
        const response = await fetch(`/api/parks/${parkId}/weather`);
        const data = await response.json();

        if (response.ok) {
            const weather = data.weather;
            document.getElementById('weatherInfo').innerHTML = `
                <div class="weather-card">
                    <h3>Current Weather</h3>
                    <div class="weather-display">
                        <img src="${weather.icon_url}" alt="${weather.description}">
                        <div class="temperature">${weather.temperature}°C</div>
                        <div class="description">${weather.description}</div>
                    </div>
                    <div class="weather-details">
                        <div>Humidity: ${weather.humidity}%</div>
                        <div>Wind: ${weather.wind_speed} km/h</div>
                        <div>Pressure: ${weather.pressure} hPa</div>
                    </div>
                </div>
            `;
        }
    } catch (error) {
        console.error('Error loading weather:', error);
    }
}
\end{lstlisting}

\section{Responsive Design}

The frontend implements responsive design principles using CSS Grid and Flexbox for optimal viewing across devices.

\begin{lstlisting}[language=CSS, caption=Responsive Layout]
.park-details {
    display: grid;
    grid-template-columns: 2fr 1fr;
    gap: 2rem;
    margin: 2rem 0;
}

.weather-card {
    background: white;
    border-radius: 8px;
    padding: 1.5rem;
    box-shadow: 0 2px 4px rgba(0,0,0,0.1);
    margin-bottom: 1rem;
}

.weather-display {
    display: flex;
    align-items: center;
    gap: 1rem;
    margin-bottom: 1rem;
}

@media (max-width: 768px) {
    .park-details {
        grid-template-columns: 1fr;
        gap: 1rem;
    }

    .weather-display {
        flex-direction: column;
        text-align: center;
    }
}
\end{lstlisting}

% API Documentation
\chapter{API Documentation}

\section{API Overview}

The Tunisia National Parks API provides RESTful endpoints for accessing and managing environmental data. All endpoints return JSON responses and follow RESTful conventions for HTTP methods and status codes.

\section{Authentication}

The API uses JWT (JSON Web Token) based authentication for protected endpoints.

\begin{lstlisting}[language=bash, caption=Authentication Example]
# Obtain access token
curl -X POST "http://localhost:8002/api/auth/token" \
  -H "Content-Type: application/json" \
  -d '{"username":"admin","password":"your-password"}'

# Use token in subsequent requests
curl -H "Authorization: Bearer YOUR_TOKEN_HERE" \
  http://localhost:8002/api/parks
\end{lstlisting}

\section{API Endpoints}

\subsection{Parks Management}

\begin{table}[H]
\centering
\caption{Parks API Endpoints}
\label{tab:parks-endpoints}
\begin{tabular}{@{}lll@{}}
\toprule
Method & Endpoint & Description \\
\midrule
GET & /api/parks & List all parks with filtering \\
GET & /api/parks/\{id\} & Get specific park details \\
POST & /api/parks & Create new park (admin) \\
PUT & /api/parks/\{id\} & Update park (admin) \\
DELETE & /api/parks/\{id\} & Delete park (admin) \\
\bottomrule
\end{tabular}
\end{table}

\subsection{Species Management}

\begin{table}[H]
\centering
\caption{Species API Endpoints}
\label{tab:species-endpoints}
\begin{tabular}{@{}lll@{}}
\toprule
Method & Endpoint & Description \\
\midrule
GET & /api/species & List all species \\
GET & /api/species/\{id\} & Get species details \\
GET & /api/parks/\{id\}/species & Species in specific park \\
POST & /api/species & Create new species (admin) \\
PUT & /api/species/\{id\} & Update species (admin) \\
DELETE & /api/species/\{id\} & Delete species (admin) \\
\bottomrule
\end{tabular}
\end{table}

\subsection{Weather Integration}

\begin{table}[H]
\centering
\caption{Weather API Endpoints}
\label{tab:weather-endpoints}
\begin{tabular}{@{}lll@{}}
\toprule
Method & Endpoint & Description \\
\midrule
GET & /api/parks/\{id\}/weather & Current weather for park \\
GET & /api/parks/\{id\}/forecast & Weather forecast \\
GET & /api/weather/current & General weather endpoint \\
\bottomrule
\end{tabular}
\end{table}

\section{Request/Response Examples}

\subsection{Parks List Response}

\begin{lstlisting}[language=JSON, caption=Parks List API Response]
{
  "data": [
    {
      "id": 1,
      "name": "Ichkeul National Park",
      "governorate": "Bizerte",
      "description": "UNESCO World Heritage wetland site",
      "latitude": 37.1617,
      "longitude": 9.6742,
      "area_km2": 12600,
      "images": [
        "/uploads/parks/1/ichkeul_park_1.jpg",
        "/uploads/parks/1/ichkeul_park_2.jpg"
      ]
    }
  ],
  "total": 17,
  "skip": 0,
  "limit": 10
}
\end{lstlisting}

\subsection{Weather Data Response}

\begin{lstlisting}[language=JSON, caption=Weather API Response]
{
  "park_id": 1,
  "park_name": "Ichkeul National Park",
  "weather": {
    "temperature": 22,
    "feels_like": 24,
    "temp_min": 18,
    "temp_max": 26,
    "humidity": 65,
    "pressure": 1013,
    "description": "partiellement nuageux",
    "icon": "02d",
    "icon_url": "https://openweathermap.org/img/wn/02d@2x.png",
    "wind_speed": 12.5,
    "wind_direction": 225,
    "clouds": 45,
    "visibility": 8.5,
    "sunrise": 1640995200,
    "sunset": 1641031200,
    "timezone": 3600,
    "city_name": "Bizerte"
  }
}
\end{lstlisting}

\section{Error Handling}

The API uses standard HTTP status codes and provides detailed error messages.

\begin{table}[H]
\centering
\caption{Common HTTP Status Codes}
\label{tab:http-codes}
\begin{tabular}{@{}lll@{}}
\toprule
Status Code & Meaning & Description \\
\midrule
200 & OK & Request successful \\
201 & Created & Resource created successfully \\
400 & Bad Request & Invalid request parameters \\
401 & Unauthorized & Authentication required \\
403 & Forbidden & Insufficient permissions \\
404 & Not Found & Resource not found \\
422 & Unprocessable Entity & Validation error \\
500 & Internal Server Error & Server error \\
\bottomrule
\end{tabular}
\end{table}

% Installation & Setup
\chapter{Installation and Setup}

\section{System Requirements}

\subsection{Minimum Requirements}
\begin{itemize}
    \item Python 3.11 or higher
    \item 4GB RAM
    \item 2GB available disk space
    \item Internet connection for external API integrations
\end{itemize}

\subsection{Recommended Requirements}
\begin{itemize}
    \item Python 3.11+
    \item 8GB RAM
    \item SSD storage for improved performance
    \item PostgreSQL database for production use
\end{itemize}

\section{Installation Steps}

\subsection{Option 1: Docker Deployment (Recommended)}

\begin{lstlisting}[language=bash, caption=Docker Installation]
# Clone the repository
git clone https://github.com/Ranim-alouii/tunisia-national-parks-api.git
cd tunisia-national-parks-api

# Start all services
docker-compose up --build

# Access the application
# Frontend: http://localhost:8000
# API Docs: http://localhost:8000/docs
\end{lstlisting}

\subsection{Option 2: Local Development Setup}

\begin{lstlisting}[language=bash, caption=Local Installation]
# Clone repository
git clone https://github.com/Ranim-alouii/tunisia-national-parks-api.git
cd tunisia-national-parks-api

# Create virtual environment
python -m venv venv
source venv/bin/activate  # Windows: venv\Scripts\activate

# Install dependencies
pip install -r requirements.txt

# Configure environment variables
cp .env.example .env
# Edit .env with your settings

# Initialize database
python -c "from database import init_db; init_db()"

# Start development server
python main.py

# Access application
# Frontend: http://localhost:8002
# API Docs: http://localhost:8002/docs
\end{lstlisting}

\section{Environment Configuration}

Create a comprehensive \texttt{.env} file for proper application configuration:

\begin{lstlisting}[language=bash, caption=Environment Configuration]
# Database Configuration
DATABASE_URL=sqlite:///./tunisia_parks.db

# Security Settings
SECRET_KEY=your-secure-secret-key-here
ADMIN_USERNAME=admin
ADMIN_PASSWORD=your-secure-admin-password

# External API Keys (Optional)
OPENWEATHER_API_KEY=your_openweather_key
UNSPLASH_ACCESS_KEY=your_unsplash_key
GOOGLE_PLACES_API_KEY=your_google_key
NEWSAPI_API_KEY=your_newsapi_key

# Application Settings
ENVIRONMENT=development
DEBUG=true
LOG_LEVEL=INFO

# CORS Settings
ALLOWED_ORIGINS=http://localhost:3000,http://localhost:8002
\end{lstlisting}

\section{Database Initialization}

\subsection{Development Database}
The application uses SQLite by default for development, which requires no additional setup.

\subsection{Production Database}
For production deployments, PostgreSQL with PostGIS extension is recommended:

\begin{lstlisting}[language=bash, caption=PostgreSQL Setup]
# Install PostgreSQL and PostGIS
sudo apt-get install postgresql-15-postgis-3

# Create database and user
createdb tunisia_parks
createuser tunisia_user
psql -c "GRANT ALL PRIVILEGES ON DATABASE tunisia_parks TO tunisia_user;"

# Enable PostGIS extension
psql tunisia_parks -c "CREATE EXTENSION postgis;"
\end{lstlisting}

\section{Testing Setup}

\begin{lstlisting}[language=bash, caption=Testing Configuration]
# Run all tests
pytest

# Run with coverage
pytest --cov=. --cov-report=html

# Run API integration tests
python test_localhost_apis.py

# Run specific test file
pytest tests/test_parks.py -v
\end{lstlisting}

% Conclusion
\chapter{Conclusion}

\section{Project Summary}

The Tunisia National Parks API represents a comprehensive solution for environmental data management and biodiversity information systems. The project successfully delivers a robust, scalable web service that integrates modern web technologies with environmental monitoring capabilities.

\section{Achievements}

\subsection{Technical Accomplishments}
\begin{itemize}
    \item Implementation of a high-performance REST API using FastAPI framework
    \item Development of comprehensive database schema supporting complex environmental data relationships
    \item Integration of external APIs for enhanced functionality (weather, imagery, news)
    \item Creation of responsive web interface with interactive mapping capabilities
    \item Implementation of JWT-based authentication and authorization systems
    \item Development of gamification features for user engagement
\end{itemize}

\subsection{Functional Achievements}
\begin{itemize}
    \item Complete coverage of 18 Tunisian national parks with detailed information
    \item Comprehensive species database with 250+ entries and conservation status
    \item User-generated content system with reviews and ratings
    \item Multi-language support for broader accessibility
    \item Real-time weather integration and forecasting
    \item Advanced search and filtering capabilities
\end{itemize}

\section{Limitations and Constraints}

\subsection{Current Limitations}
\begin{itemize}
    \item Dependency on external API services for enhanced features
    \item Limited offline functionality without internet connectivity
    \item Scalability constraints with SQLite for large-scale deployments
    \item Manual data validation requirements for user-generated content
\end{itemize}

\section{Future Enhancements}

\subsection{Short-term Improvements}
\begin{itemize}
    \item Implementation of automated testing pipelines
    \item Enhanced mobile application development
    \item Integration with IoT environmental sensors
    \item Advanced analytics and reporting dashboards
\end{itemize}

\subsection{Long-term Vision}
\begin{itemize}
    \item Development of machine learning models for species identification
    \item Integration with satellite imagery for habitat monitoring
    \item Expansion to other North African countries
    \item Implementation of blockchain-based conservation certificates
\end{itemize}

\section{Impact and Significance}

The Tunisia National Parks API contributes significantly to environmental conservation efforts by:

\begin{itemize}
    \item Providing accessible environmental data for research and policy making
    \item Promoting sustainable tourism through informed decision making
    \item Supporting biodiversity conservation through data-driven approaches
    \item Enabling citizen science participation in environmental monitoring
    \item Facilitating international collaboration in conservation initiatives
\end{itemize}

\section{Final Remarks}

This project demonstrates the effective application of modern web technologies to address environmental challenges. The combination of robust backend architecture, comprehensive data management, and user-friendly interfaces creates a platform that serves both technical users and general public, contributing to the broader goals of environmental sustainability and biodiversity conservation in Tunisia.

\end{document}
